\documentclass[11pt,a4paper]{article}
\usepackage{kotex}
\usepackage{fontspec}
\setmainfont{Noto Sans CJK KR}
\usepackage[a4paper,margin=1in]{geometry}
\usepackage{amsmath,amssymb,mathtools,bm}
\usepackage{booktabs,tabularx,array,longtable,multirow}
\usepackage{graphicx}
\usepackage{hyperref}
\usepackage{xcolor}
\usepackage{caption}
\usepackage{enumitem}
\usepackage{float}
\usepackage{setspace}
\usepackage{microtype}
\usepackage{fancyhdr}
\usepackage{titlesec}
\usepackage{url}
\graphicspath{{./}{/mnt/data/}}

\hypersetup{
  colorlinks=true,
  linkcolor=blue!50!black,
  urlcolor=blue!50!black,
  citecolor=blue!50!black
}
\Urlmuskip=0mu plus 1mu
\setstretch{1.12}
\setlength{\parskip}{0.45em}
\setlength{\parindent}{0pt}
\pagestyle{fancy}
\setlength{\headheight}{14pt}
\fancyhf{}
\fancyhead[L]{워킹노트 초안}
\fancyhead[R]{\thepage}
\titleformat{\section}{\large\bfseries}{\thesection.}{0.5em}{}
\titleformat{\subsection}{\normalsize\bfseries}{\thesubsection.}{0.5em}{}

\newcolumntype{Y}{>{\raggedright\arraybackslash}X}
\newcommand{\Prob}{\mathbb{P}}
\newcommand{\logit}{\operatorname{logit}}
\newcommand{\HHI}{\operatorname{HHI}}

\title{\vspace{-1.0em}\textbf{상장사 비트코인 보유비중, 보유분포 HHI, 공개 계약표기 프록시를 이용한 보수적 ordered-stage 추정}}
\author{Theresno2ndbtc}
\date{2026년 3월}

\begin{document}
\maketitle
\vspace{-1.4em}

\begin{abstract}
본 노트는 공개 상장사 비트코인 보유비중과 보유분포의 Herfindahl--Hirschman Index(HHI), 그리고 공개 BTC-denominated 자금조달 프록시와 merchant support 프록시를 결합하여 broad public-denomination stage의 확률을 ordered-stage 틀에서 추정한다. 설명변수는 (i) 공개 상장사 전체 BTC 보유량의 총공급 2,100만 대비 비중, (ii) 공개 상장사 보유분포의 HHI, (iii) 공개 BTC-denominated 자금조달 breadth와 merchant support를 결합한 지원지수다. 종속변수는 treasury-only, financing-pilot, public-financing-cluster, broad public-denomination의 네 단계다. 마지막 단계는 공개 BTC-denominated loan/bond 사례뿐 아니라, 반복적 운영계약(invoice/payroll/lease/subscription)의 공적 확인을 요구하도록 엄격하게 정의한다. 2025년 1월부터 2026년 3월 11일까지의 월말 표본에서 broad public-denomination stage는 아직 관측되지 않는다. 따라서 최상위 절단점은 엄격한 단계정의를 반영하는 prior와 함께 MAP ordered logit으로 규율된다. 기준 ordered logit 추정에서 보유비중 계수는 양(+), HHI 계수는 음(-), 지원지수 계수는 양(+)으로 나타난다. 2026년 3월 11일 현재 관측치($x_t=5.50\%$, HHI$=0.4154$)를 대입하면 broad public-denomination stage 확률은 ordered logit 기준 11.2\%, ordered probit 민감도 기준 3.1\%다. 현재 HHI와 지원지수를 고정할 때 Stage 3 확률이 50\%, 70\%, 90\%에 도달하는 보유비중 임계치는 ordered logit 기준 각각 8.78\%, 10.11\%, 12.24\%로 계산된다. 본 결과는 $x$ 의 상승이 단순한 treasury adoption에 그치지 않고, 일정 수준을 넘어서면 BTC-denominated 계약과 가격표시의 확산을 통해 UoA 전환 확률을 높이는 방향으로 작동할 수 있음을 시사한다.
\end{abstract}

\section{문제 설정}

경제 전체의 unit-of-account(UoA) 채택은 ``무엇으로 가격, 임금, 청구서, 대출, 채권을 적느냐''의 문제다. 그러나 이 변수는 공식 통계로 직접 관측되기 어렵다. 공개자료에서 연속적으로 집계 가능한 것은 주로 세 가지다. 첫째, 공개 상장사의 BTC 보유량. 둘째, 그 보유가 소수 기업에 집중되어 있는지 여부. 셋째, 공개 BTC-denominated 기업금융 사례와 merchant 네트워크의 범위다.

본 노트는 broad UoA를 곧바로 측정하지 않고, 공개자료로 확인 가능한 단계적 전개를 다음 네 수준으로 정리한다.

\begin{table}[H]
\centering
\caption{단계 정의}
\label{tab:stagedef}
\begin{tabularx}{\textwidth}{@{}lY@{}}
\toprule
단계 & 정의 \\
\midrule
Stage 0 & Treasury-only 단계. 공개 상장사 보유가 증가하더라도 공개 BTC-denominated 자금조달 사례가 아직 없다. \\
Stage 1 & Financing-pilot 단계. 공개 BTC-denominated loan/bond/convertible 사례가 산발적으로 등장하지만 breadth가 좁다. \\
Stage 2 & Public-financing-cluster 단계. 공개 BTC-denominated 자금조달 사례가 다수 누적되고 merchant support도 확대되지만, 반복적 운영계약(invoice/payroll/lease/subscription)의 공적 확인은 아직 없다. \\
Stage 3 & Broad public-denomination 단계. Stage 2 조건에 더해 반복적 운영계약의 공개 BTC-denomination이 확인되는 구간이다. 본 노트의 표본에서는 이 단계가 아직 관측되지 않는다. \\
\bottomrule
\end{tabularx}
\end{table}

이 정의에서 Stage 3는 경제 전체의 완성된 UoA와 동일하지는 않지만, 적어도 공개 계약표기 차원에서의 상위 단계를 뜻한다. 따라서 Stage 3 확률이 높아지려면 단순한 treasury 보유 증가만으로는 부족하고, 보유의 분산, 공개 자금조달 breadth, merchant support, 운영계약의 반복성까지 함께 갖추어져야 한다.

\section{데이터}

\subsection{공개 상장사 보유 BTC 비중}

시점 $t$의 보유비중은 다음과 같이 정의한다.
\begin{equation}
 x_t = 100 \times \frac{B_t^{\mathrm{public}}}{21{,}000{,}000}.
\end{equation}
여기서 $B_t^{\mathrm{public}}$는 공개 상장사의 총 BTC 보유량이다. 본 노트의 최신 관측치는 2026년 3월 11일 1,156,000 BTC로 두었고, 이는 총공급 대비 5.50\%다.

\begin{table}[H]
\centering
\caption{공개 상장사 보유 BTC 앵커 시계열}
\label{tab:anchors_pub2}
\begin{tabularx}{\textwidth}{@{}lrrY@{}}
\toprule
날짜 & 공개 상장사 보유 BTC & 총공급 대비 비중(\%) & 비고 \\
\midrule
2024-12-31 & 590,649 & 2.81 & 2024년말 공개 상장사 총보유량 \\
2025-03-31 & 688,000 & 3.28 & 2025년 1분기말 \\
2025-06-30 & 847,000 & 4.03 & 2025년 2분기말 \\
2025-09-30 & 1,022,700 & 4.87 & 2025년 3분기말 \\
2025-12-31 & 1,100,000 & 5.24 & 2025년말 \\
2026-03-11 & 1,156,000 & 5.50 & 2026년 3월 11일 현재 \\
\bottomrule
\end{tabularx}
\end{table}

월별 표본은 위 앵커 관측치 사이를 시간기준 선형보간하여 구성했다.

\subsection{공개 상장사 보유분포의 HHI}

시점 $t$의 집중도는 다음과 같이 정의한다.
\begin{equation}
 h_t = \HHI_t = \sum_{i=1}^{N_t} s_{it}^2, \qquad s_{it}=\frac{B_{it}}{B_t^{\mathrm{public}}}.
\end{equation}
여기서 $B_{it}$는 시점 $t$에 $i$번째 공개 상장사가 보유한 BTC이며, $N_t$는 추적 가능한 공개 상장사 수다.

2026년 3월 11일 현재 cross-section에서는 공개 상장사 총보유량이 1,156,000 BTC, Strategy 보유량이 738,731 BTC이므로 top-1 share는 63.90\%다. 같은 날짜의 공개 보유분포에서 식별 가능한 보유량과 residual tail을 결합하면 현재 HHI 프록시는
\begin{equation}
 \HHI_{2026-03-11} \approx 0.4154
\end{equation}
로 계산된다. 10,000 스케일로는 약 4,154다.

월별 HHI 시계열은 Strategy 비중의 월별 경로와 현재 시점의 residual concentration shape를 결합해 구성했다.

\subsection{공개 자금조달 breadth와 merchant support}

공개 BTC-denominated 기업금융 breadth는 일곱 개 공적 사례로 구성한다. 본 노트는 2025년 3월부터 2025년 12월 사이 공개된 BTC-denominated convertible bond, BTC loan, BTC bond 사례를 event panel로 정리한다. 각 월말에 활성화된 누적 비율을
\begin{equation}
 f_t = \frac{k_t}{7},
\end{equation}
로 둔다. 여기서 $k_t$는 시점 $t$까지 공적으로 확인된 사례 수다.

merchant support는 BTC Map의 공개 merchant 수 시계열을 월별로 보간한 뒤,
\begin{equation}
 m_t = \frac{\text{BTC-accepting places}_t}{30{,}000}
\end{equation}
로 정규화했다. Stage 3의 핵심은 직접 가격표기와 반복계약이지만, 공개 월별 시계열이 충분하지 않기 때문에 merchant support는 \emph{보조적 인프라 프록시}로 사용한다.

두 변수를 다음과 같이 결합한다.
\begin{equation}
 g_t = 0.65 f_t + 0.35 m_t.
\end{equation}
가중치 0.65/0.35는 공개 BTC-denominated financing breadth가 merchant acceptance보다 더 직접적인 계약표기 신호라는 판단을 반영한다.

\subsection{엄격한 ordered-stage 종속변수}

본 노트의 관측 stage는 다음의 규칙으로 정의한다.
\begin{equation}
S_t=
\begin{cases}
0, & f_t=0\ \text{ and } x_t<3.4,\\
1, & f_t<0.5\ \text{ or } m_t<0.55,\\
2, & f_t\ge 0.5,\ m_t\ge 0.55,\ \text{and 반복적 운영계약의 공적 확인 없음},\\
3, & \text{반복적 운영계약의 공적 확인 존재}.
\end{cases}
\end{equation}
표본기간에는 반복적 BTC-denominated invoice/payroll/lease/subscription의 공적 확인이 없으므로 Stage 3는 아직 관측되지 않는다. 이 점이 본 노트의 핵심이다. 현재 공개 상장사 보유비중이 5.5\% 수준이라 해도, strict stage 정의 아래에서는 여전히 broad public-denomination의 초기 전 단계로 해석할 수 있다.

\section{모형}

잠재 연속변수 $S_t^*$를
\begin{equation}
S_t^* = \beta_x \tilde x_t + \beta_h \tilde h_t + \beta_g \tilde g_t + \varepsilon_t,
\end{equation}
로 둔다. 여기서 $\tilde x_t, \tilde h_t, \tilde g_t$는 표준화된 보유비중, HHI, 지원지수다. ordered logit 기준에서 $\varepsilon_t$는 logistic 분포를 따른다.

관측된 stage는 절단점 $(\kappa_1,\kappa_2,\kappa_3)$에 의해
\begin{equation}
S_t=
\begin{cases}
0, & S_t^* \le \kappa_1, \\
1, & \kappa_1 < S_t^* \le \kappa_2, \\
2, & \kappa_2 < S_t^* \le \kappa_3, \\
3, & S_t^* > \kappa_3
\end{cases}
\end{equation}
로 결정된다.

따라서 ordered logit의 단계확률은
\begin{align}
\Prob(S_t=0\mid X_t) &= \Lambda(\kappa_1-\eta_t), \\
\Prob(S_t=1\mid X_t) &= \Lambda(\kappa_2-\eta_t)-\Lambda(\kappa_1-\eta_t), \\
\Prob(S_t=2\mid X_t) &= \Lambda(\kappa_3-\eta_t)-\Lambda(\kappa_2-\eta_t), \\
\Prob(S_t=3\mid X_t) &= 1-\Lambda(\kappa_3-\eta_t),
\end{align}
여기서 $\eta_t = \beta_x \tilde x_t + \beta_h \tilde h_t + \beta_g \tilde g_t$다.

추정은 약한 정규 사전분포를 둔 MAP 방식으로 한다.
\begin{equation}
\beta_x,\beta_h,\beta_g \sim \mathcal{N}(0,1.2^2),
\end{equation}
절단점은
\begin{equation}
\kappa_1 = a_1, \qquad \kappa_2 = a_1 + e^{d_2}, \qquad \kappa_3 = a_1 + e^{d_2} + e^{d_3}
\end{equation}
로 parameterize했고,
\begin{equation}
 a_1 \sim \mathcal{N}(-1,1^2),\qquad d_2 \sim \mathcal{N}(\log 1.4,0.25^2),\qquad d_3 \sim \mathcal{N}(\log 3.8,0.16^2)
\end{equation}
를 두었다. 마지막 prior는 Stage 3가 반복적 운영계약의 공적 확인을 요구하는 엄격한 단계라는 사실을 반영한다. 표본에서 Stage 3가 아직 관측되지 않기 때문에, 최상위 절단점은 데이터와 단계정의를 동시에 반영하는 prior-disciplined parameter가 된다.

ordered probit도 민감도 점검으로 함께 계산했다.

\section{추정 결과}

\subsection{기준 ordered logit}

표 \ref{tab:mainres_strict}는 기준 ordered logit 추정치를 보여준다.

\begin{table}[H]
\centering
\caption{기준 ordered logit 추정치}
\label{tab:mainres_strict}
\begin{tabularx}{0.88\textwidth}{@{}lrrY@{}}
\toprule
모수 & 추정치 & 부호 & 해석 \\
\midrule
$\beta_x$ & 0.548 & + & 보유비중이 높을수록 더 높은 stage로 이동 \\
$\beta_h$ & -0.637 & -- & 동일한 보유비중에서도 HHI가 높을수록 stage 상승이 지연 \\
$\beta_g$ & 0.498 & + & 공개 financing breadth와 merchant support가 높을수록 stage 상승 \\
$\kappa_1$ & -1.907 &  & Stage 0/1 경계 \\
$\kappa_2$ & -0.344 &  & Stage 1/2 경계 \\
$\kappa_3$ & 3.822 &  & Stage 2/3 경계 \\
\bottomrule
\end{tabularx}
\end{table}

계수의 부호는 직관과 부합한다. $\beta_x>0$는 공개 상장사 BTC 보유비중이 높을수록 더 높은 단계로 이동할 가능성이 커짐을 뜻한다. 반면 $\beta_h<0$는 동일한 보유비중에서도 보유분포의 HHI가 높을수록 상위 단계로의 이동이 늦어진다는 뜻이다. $\beta_g>0$는 공개 financing breadth와 merchant support가 stage 상승에 기여함을 의미한다.

\subsection{현재 시점의 단계확률}

2026년 3월 11일 현재 관측치는 $x_t=5.50\%$, $\HHI_t=0.4154$, $g_t=0.936$이다. 이를 대입하면 기준 ordered logit의 단계확률은 표 \ref{tab:currprobs_strict}와 같다.

\begin{table}[H]
\centering
\caption{2026년 3월 11일 현재 단계확률}
\label{tab:currprobs_strict}
\begin{tabular}{@{}lrr@{}}
\toprule
Stage & 의미 & ordered logit 확률 \\
\midrule
0 & treasury-only 이전 구간 & 2.5\% \\
1 & financing-pilot 단계 & 8.5\% \\
2 & public-financing-cluster 단계 & 77.9\% \\
3 & broad public-denomination 단계 & 11.2\% \\
\bottomrule
\end{tabular}
\end{table}

즉 현재 관측치는 Stage 2가 지배적이며, Stage 3는 아직 low-probability tail에 머문다. 이는 ``공개 treasury와 공개 financing은 상당히 진행되었지만, 반복적 운영계약 차원의 계약표기는 여전히 이르다''는 서술과 양립한다.

\subsection{현재 조건에서의 Stage 3 임계치}

현재 HHI와 지원지수($h_t=0.4154$, $g_t=0.936$)를 고정할 때, Stage 3 확률이 특정 수준에 도달하는 보유비중 $x^*(p)$는
\begin{equation}
 \Prob(S_t=3\mid x_t,h_t,g_t)=p
\end{equation}
를 만족하는 $x_t$로 정의한다. ordered logit 기준 임계치는 표 \ref{tab:thresholds_strict}와 같다.

\begin{table}[H]
\centering
\caption{현재 조건 기준 Stage 3 임계치}
\label{tab:thresholds_strict}
\begin{tabular}{@{}lrr@{}}
\toprule
목표 확률 $p$ & 임계치 $x^*(p)$ (\% of 21M) & BTC 수량 환산 \\
\midrule
50\% & 8.78\% & 1,843,366 BTC \\
70\% & 10.11\% & 2,124,070 BTC \\
90\% & 12.24\% & 2,571,292 BTC \\
\bottomrule
\end{tabular}
\end{table}

이 숫자는 Stage 3를 엄격하게 정의했기 때문에, 단순 treasury adoption 임계치보다 훨씬 높다. 현재 5.50\%는 Stage 3의 50\% 확률 임계치인 8.78\%보다 유의하게 낮다.

\begin{figure}[H]
\centering
\includegraphics[width=0.88\textwidth]{btc_strict_stage3_probability_curve.png}
\caption{현재 HHI와 지원지수를 고정했을 때의 Stage 3 확률 곡선}
\label{fig:curve}
\end{figure}

\subsection{ordered probit 민감도 점검}

링크함수를 normal cdf로 바꾼 ordered probit 민감도 점검에서는 현재 Stage 3 확률이 3.1\%로 더 낮게 나온다. 같은 방식으로 계산한 임계치는 표 \ref{tab:probit_sens}와 같다.

\begin{table}[H]
\centering
\caption{ordered probit 민감도 점검: Stage 3 임계치}
\label{tab:probit_sens}
\begin{tabular}{@{}lrr@{}}
\toprule
목표 확률 $p$ & 임계치 $x^*(p)$ (\% of 21M) & BTC 수량 환산 \\
\midrule
50\% & 8.59\% & 1,803,836 BTC \\
70\% & 9.46\% & 1,985,832 BTC \\
90\% & 10.71\% & 2,248,605 BTC \\
\bottomrule
\end{tabular}
\end{table}

ordered logit과 ordered probit은 절대수준에서 차이를 보이지만, 두 경우 모두 현재 5.50\%를 broad public-denomination 단계의 고확률 구간으로 보지는 않는다. 즉 결론의 방향은 동일하다.

\section{해석}

이 노트의 핵심 메시지는 다음과 같다.

첫째, 공개 상장사 BTC 보유비중은 여전히 중요한 설명변수다. 보유비중이 높아질수록 더 높은 단계로 이동할 잠재력은 커진다.

둘째, 그러나 HHI가 높으면 같은 보유비중에서도 상위 단계로의 이동이 늦어진다. 현재 공개 상장사 보유는 여전히 소수 기업에 상당히 집중되어 있다.

셋째, broad public-denomination 단계는 단순한 BTC treasury 보유나 공개 financing 사례만으로 열리지 않는다. 반복적 운영계약의 공적 확인이 없으면 Stage 3는 높은 절단점을 갖는 것이 자연스럽다.

넷째, 따라서 현재 $x_t\approx 5.5\%$라는 관측치만으로 ``UoA가 이미 고확률 구간에 진입했다''고 읽는 것은 과도하다. 본 노트의 strict stage 정의와 conservative calibration 아래에서는 현재 상태를 public-financing-cluster 단계가 지배적이지만 broad public-denomination은 아직 초기라고 해석하는 것이 더 자연스럽다.

\section{한계}

첫째, Stage 3는 표본에서 아직 관측되지 않는다. 따라서 최상위 절단점은 순수한 빈도주의 식별이 아니라, 단계정의를 반영한 prior와 함께 규율된다.

둘째, merchant support는 가격표기 그 자체가 아니라 보조적 인프라 프록시다. 향후 sats 직접 가격표시 merchant 비중, BTC invoice/payroll/lease/subscription 데이터가 축적되면 $g_t$를 더 직접적인 UoA 프록시로 대체할 수 있다.

셋째, 공개 BTC-denominated financing event panel은 탐색적 표본이다. 그러나 본 노트의 목적은 경제 전체 UoA를 직접 계측하기보다, 공개자료에 의해 확인되는 \emph{엄격한 public-denomination 단계}의 확률을 보수적으로 요약하는 데 있다.

\section{결론}

현재 공개 상장사 BTC 보유비중은 약 5.50\%까지 상승했지만, 보유 집중도와 운영계약 차원의 공적 확인 부족을 함께 고려하면 broad public-denomination 단계의 확률은 아직 높지 않다. 기준 ordered logit에서는 11.2\%, ordered probit 민감도 점검에서는 3.1\%이며, Stage 3의 50\% 임계치는 대략 8.6\%--8.8\%, 70\% 임계치는 9.5\%--10.1\%, 90\% 임계치는 10.7\%--12.2\% 범위로 계산된다. 따라서 현재의 공개 데이터는 ``treasury와 공개 financing은 상당히 진행되었으나, broad public-denomination 단계는 아직 이르다''는 해석과 더 잘 부합한다. 다만 본 결과는 $x$ 의 상승이 단순한 treasury adoption에 그치지 않고, 일정 수준을 넘어서면 BTC-denominated 계약과 가격표시의 확산을 통해 UoA 전환 확률을 높이는 방향으로 작동할 수 있음을 보여준다. 다시 말해 현 시점은 아직 초기 단계로 해석하는 것이 적절하지만, 공개 상장사 보유비중의 추가 상승은 공개 계약표기와 운영단계 사용의 누적 확산을 뒷받침하면서 향후 UoA 형성 가능성을 높이는 방향의 조건 변화로 읽을 수 있다.

\section*{부록 A. 데이터 메모}

표본은 2025년 1월부터 2026년 3월 11일까지 15개 월말 관측치다. 공개 상장사 총보유 BTC와 Strategy 보유량은 공개 treasury tracker를 기초로 구성했고, HHI는 현재 cross-section의 보유분포와 residual tail 가정을 이용해 프록시화했다. BTC-denominated financing breadth는 공개 기업대출/전환사채/회사채 사례 일곱 개를 누적 카운트하여 구성했고, merchant support는 BTC Map 공개 merchant 수의 월별 보간치를 사용했다.

\end{document}
